% ---------------------------------------------------------------------------
% Rysunki TikZ manuala Stardust NEXT
% ---------------------------------------------------------------------------
\tikzset{
  box/.style={draw=sdnavy,fill=white,rounded corners=2pt,align=center,
    font=\sffamily\scriptsize,inner sep=3pt,minimum height=7mm},
  boxo/.style={box,fill=sdorange!12,draw=sdorange},
  boxg/.style={box,fill=sdlight,draw=sdgray},
  lbl/.style={font=\sffamily\scriptsize,align=left},
  lblc/.style={font=\sffamily\scriptsize,align=center},
  radio/.style={-{Stealth[length=2mm]},decorate,
    decoration={snake,amplitude=0.6mm,segment length=2.2mm,post length=2mm},draw=sdorange,thick},
  net/.style={-{Stealth[length=2mm]},draw=sdnavy,thick},
  dim/.style={{Stealth[length=1.5mm]}-{Stealth[length=1.5mm]},draw=sdgray},
}

% ---------------------------------------------------------------- zestaw lotny
\newcommand{\rysZestaw}{%
\begin{tikzpicture}[x=1cm,y=1cm,font=\sffamily\scriptsize]
  % powłoka
  \shade[ball color=white,draw=sdgray] (0,11.2) ellipse (1.25 and 1.45);
  \draw[sdgray] (-0.18,9.85) -- (-0.12,9.55) -- (0.12,9.55) -- (0.18,9.85);
  \node[lbl,anchor=west] at (1.6,11.2) {\textbf{Powłoka} lateksowa z helem\\średnica przy starcie ok. 1,5–2 m};
  \node[lbl,anchor=west] at (1.6,9.7) {\textbf{Szyjka} zamknięta: skręcona, związana,\\zagięta w górę i owinięta taśmą};
  % krętlik
  \draw[thick] (0,9.55) -- (0,9.2);
  \draw[thick,fill=sdlight] (0,9.05) circle (0.15);
  \node[lbl,anchor=west] at (1.6,9.0) {\textbf{Krętlik} i krótka pętla mocująca};
  % linka do spadochronu
  \draw[thick] (0,8.9) -- (0,8.1);
  % spadochron
  \fill[sdorange] (-0.6,7.3) -- (0,8.1) -- (0.6,7.3) -- cycle;
  \draw[sdgray] (-0.6,7.3) -- (0,6.4) -- (0.6,7.3);
  \node[lbl,anchor=west] at (1.6,7.5) {\textbf{Spadochron} (stale otwarty)};
  % linka główna i zapas
  \draw[thick] (0,6.4) -- (0,2.4);
  \draw[sdorange,thick,densely dashed] (0,9.2) .. controls (-1.6,8.4) and (-1.3,4.2) .. (-0.35,2.4);
  \node[lbl,anchor=east,text=sdorange] at (-1.5,6.3) {Linka\\odcinacza};
  \draw[dim] (0.35,6.35) -- node[right,lbl]{\textbf{zapas linki} 10–15\,m\\(miejsce na otwarcie spadochronu\\mimo skręconych linek)} (0.35,2.45);
  % gondola platformy
  \draw[fill=white,draw=sdnavy,thick] (-0.75,1.6) rectangle (0.75,2.4);
  \draw[fill=sdorange] (-0.45,2.4) rectangle (-0.15,2.55);
  \node[lblc] at (0,2.0) {platforma};
  \node[lbl,anchor=west] at (1.6,2.15) {\textbf{Gondola platformy}: tracker, OBC,\\zwalniacz odcinacza (na pokrywie), baterie};
  % gondola ladunku
  \draw[fill=white,draw=sdnavy,thick] (-0.85,0.4) rectangle (0.85,1.6);
  \node[lblc] at (0,1.0) {ładunek};
  \node[lbl,anchor=west] at (1.6,0.9) {\textbf{Gondola ładunku} (eurobox), spięta\\z platformą pasami};
  % anteny
  \draw[thick] (-0.5,0.4) -- (-0.5,-0.25);
  \draw[thick] (0.5,0.4) -- (0.5,-0.15);
  \node[lbl,anchor=west] at (1.6,-0.1) {Anteny skierowane w dół};
  % pasy
  \draw[sdorange,line width=1.2pt] (-0.1,0.4) -- (-0.1,2.4);
\end{tikzpicture}}

% ---------------------------------------------------------------- tor radiowy
\newcommand{\rysTorRadiowy}{%
\begin{tikzpicture}[node distance=6mm and 8mm,font=\sffamily\scriptsize]
  % segment lotny
  \node[boxo,minimum width=30mm] (gps) {Odbiornik GPS\\(tryb \emph{airborne})};
  \node[boxo,minimum width=30mm,below=of gps] (hab) {HABduino\\(komputer bezpieczeństwa)};
  \node[boxo,minimum width=30mm,below=of hab] (obc) {OBC\\(funkcje zapasowe)};
  \node[boxg,minimum width=30mm,below=of obc] (spot) {SPOT Trace\\(po lądowaniu)};
  \node[boxg,minimum width=30mm,below=of spot] (tag) {AirTag / SmartTag ×2\\(po lądowaniu)};
  \draw[net] (gps) -- (hab);
  \draw[net] (hab) -- (obc);
  \begin{scope}[on background layer]
    \node[draw=sdorange,dashed,rounded corners,fit=(gps)(tag),inner sep=4mm,
      label={[font=\sffamily\bfseries\scriptsize,text=sdorange]above:GONDOLA}] (gond) {};
  \end{scope}
  % stacje
  \node[box,right=34mm of hab,minimum width=34mm] (mob) {Stacja mobilna\\radiotelefon APRS (FT5DE)\\lub RTL-SDR + LNA + filtr + laptop};
  \node[box,above=7mm of mob,minimum width=34mm] (fix) {Stacja stacjonarna\\antena Yagi / dookólna, iGate};
  \node[box,below=7mm of mob,minimum width=34mm] (sieci) {Sieć APRS\\digipeatery, iGate innych\\krótkofalowców};
  \node[boxg,below=14mm of sieci,minimum width=34mm] (sat) {Satelity Globalstar};
  \node[boxg,below=6mm of sat,minimum width=34mm] (phone) {Smartfony w pobliżu\\(sieć Find My / Find Hub)};
  % chmura
  \node[box,right=12mm of mob,minimum width=26mm,fill=sdlight] (inet) {Internet\\(LTE w samochodzie)};
  \node[box,right=8mm of inet,minimum width=24mm] (sh) {SondeHub\\amateur.sondehub.org};
  \node[box,above=5mm of sh,minimum width=24mm] (aprsfi) {APRS-IS\\aprs.fi};
  \node[box,below=5mm of sh,minimum width=24mm] (spotweb) {SPOT /\\Find My (aplikacja)};
  % łącza radiowe
  \draw[radio] (hab.east) -- node[above,sloped,lblc]{RTTY 434 MHz} (fix.west);
  \draw[radio] (hab.east) -- node[above,sloped,lblc]{APRS 144,8 MHz} (mob.west);
  \draw[radio] (hab.east) -- (sieci.west);
  \draw[radio] (spot.east) -- node[below,sloped,lblc]{satelitarnie} (sat.west);
  \draw[radio] (tag.east) -- node[below,sloped,lblc]{Bluetooth} (phone.west);
  % sieć
  \draw[net] (fix.east) -| (inet.north);
  \draw[net] (mob.east) -- (inet.west);
  \draw[net] (sieci.east) -| (inet.south);
  \draw[net] (inet.east) -- (sh.west);
  \draw[net] (inet.north east) |- (aprsfi.west);
  \draw[net] (sat.east) -| (spotweb.south);
  \draw[net] (phone.east) -| ([xshift=3mm]spotweb.south);
  \node[lbl,below=1mm of sh,anchor=north west,xshift=-12mm,yshift=-13mm] {};
\end{tikzpicture}}

% ---------------------------------------------------------------- stanowisko startowe (warstwy)
\newcommand{\rysStanowisko}{%
\begin{tikzpicture}[font=\sffamily\footnotesize,x=0.62cm,y=1cm]
  % ziemia
  \fill[pattern=north east lines,pattern color=sdgray!60] (-0.6,-0.35) rectangle (10.6,0);
  \draw[thick,sdgray] (-0.6,0) -- (10.6,0);
  % warstwy (grubość przesadzona dla czytelności)
  \fill[blue!55] (0,0) rectangle (10,0.12);
  \fill[brown!60] (0.3,0.12) rectangle (9.7,0.32);
  \draw[sdgray,line width=0.8pt] (0.1,0.36) -- (9.9,0.36);
  \draw[fill=white,draw=sdgray] (1,0.4) .. controls (2,1.4) and (6,1.5) .. (8,0.55) -- (8,0.4) -- cycle;
  \fill[blue!60] (8,0.43) rectangle (9.6,0.55);
  \draw[sdgray,line width=1.3pt] (0,0.42) .. controls (1.5,1.95) and (6.5,2.05) .. (10,0.42);
  % skrzynki na rogach
  \foreach \x in {-0.55,9.65}{\draw[fill=sdnavy] (\x,0.12) rectangle ++(0.9,0.55);}
  % etykiety
  \def\opis##1##2##3##4{\draw[sdgray] (##1,##2) -- (11.0,##3) node[lbl,anchor=west]{##4};}
  \opis{9.0}{0.06}{-0.2}{1. plandeka -- ochrona przed wilgocią}
  \opis{9.2}{0.22}{0.35}{2. koc -- ochrona przed korzeniami i kamieniami}
  \opis{9.4}{0.36}{0.9}{3. folia malarska dolna (cienka)}
  \opis{8.8}{0.49}{1.45}{4. rura do tankowania w szyjce powłoki}
  \opis{5.0}{1.0}{2.0}{5. powłoka}
  \opis{7.5}{1.4}{2.55}{6. folia górna (grubsza) -- przytrzymuje powłokę}
  \opis{10.0}{0.6}{3.1}{skrzynki transportowe dociskają rogi}
  % wiatr
  \draw[-{Stealth},sdorange,thick] (2.5,2.5) -- (5.0,2.5) node[right,text=sdorange]{wiatr};
\end{tikzpicture}}

% ---------------------------------------------------------------- tankowanie
\newcommand{\rysTankowanie}{%
\begin{tikzpicture}[font=\sffamily\scriptsize,x=1cm,y=1cm]
  % butla
  \draw[fill=sdgray!40,rounded corners=4pt] (0,0) rectangle (0.9,3.6);
  \draw[fill=sdgray!60] (0.3,3.6) rectangle (0.6,3.95);
  \node[lblc,rotate=90] at (0.45,1.8) {butla z helem};
  % regulator
  \draw[fill=yellow!50!brown] (0.3,3.95) rectangle (0.95,4.2);
  \draw[fill=white,thick] (0.9,4.75) circle (0.32); \node at (0.9,4.75) {\tiny bar};
  \draw[fill=white,thick] (1.55,4.4) circle (0.28); \node at (1.55,4.4) {\tiny l/min};
  \draw (0.9,4.43) -- (0.75,4.2); \draw (1.3,4.3) -- (0.95,4.1);
  \node[lbl,anchor=south] at (0.6,5.15) {\textbf{reduktor}: manometr ciśnienia butli\\i wskaźnik przepływu};
  % wąż
  \draw[red!70!black,line width=1.5pt] (1.8,4.25) .. controls (3,4.1) and (3,1.2) .. (4.2,1.2);
  \node[lbl] at (3.3,3.0) {wąż};
  % rura
  \draw[fill=blue!60,draw=blue!80] (4.2,1.08) rectangle (5.6,1.32);
  \node[lbl,anchor=north] at (4.9,1.0) {rura do tankowania};
  % szyjka i powłoka
  \draw[fill=white,draw=sdgray] (5.6,1.05) -- (6.1,0.95) -- (6.1,1.45) -- (5.6,1.35) -- cycle;
  \shade[ball color=white,draw=sdgray] (7.9,2.0) ellipse (1.9 and 1.6);
  \draw[sdorange,thick] (5.85,1.0) -- (5.85,1.4);
  \draw[sdorange] (5.85,1.42) -- (4.6,2.5) node[lbl,anchor=south,text=sdorange]{taśma: wstępne\\uszczelnienie};
  % obciążnik
  \draw[thick] (5.95,0.95) -- (5.95,0.25);
  \draw[fill=cyan!20,draw=cyan!60!black] (5.65,-0.55) rectangle (6.25,0.25);
  \node[lbl,anchor=west] at (6.4,-0.15) {\textbf{obciążnik} (butelka z wodą)\\masa = siła nośna szyjki};
  % folia
  \draw[sdgray,line width=1pt] (5.3,3.9) .. controls (7,4.1) and (9,3.9) .. (10.2,2.4);
  \node[lbl,anchor=south] at (8.6,4.0) {folia górna trzymana przez pomocników};
  % ziemia
  \draw[thick,sdgray] (-0.3,-0.65) -- (10.5,-0.65);
\end{tikzpicture}}

% ---------------------------------------------------------------- gondola (przekrój)
\newcommand{\rysGondola}{%
\begin{tikzpicture}[font=\sffamily\scriptsize,x=1cm,y=1cm]
  % pudełko ładunku
  \draw[thick,fill=white] (0,0) rectangle (5,2.6);
  \draw[thick] (0.25,0.25) rectangle (4.75,2.35);
  \node[lblc,fill=white] at (2.5,1.3) {GONDOLA ŁADUNKU\\(eurobox, zespół badawczy)};
  % pudełko platformy
  \draw[thick,fill=white] (0,2.6) rectangle (5,4.4);
  \draw[thick] (0.25,2.75) rectangle (4.75,4.15);
  % pokrywa
  \draw[thick,fill=white] (-0.1,4.4) rectangle (5.1,4.7);
  % płyta montażowa wewnątrz
  \draw[fill=black!80] (0.3,2.78) rectangle (4.7,2.92);
  \node[lbl,anchor=west] at (5.4,2.85) {\textbf{wkładka}: płyta z plexi / sklejki / włókna\\węglowego na dnie pudełka};
  % elementy
  \draw[fill=sdorange!70] (0.5,2.92) rectangle (1.7,3.6); \node[lblc] at (1.1,3.26) {bateria};
  \draw[fill=green!40!black!40] (1.9,2.92) rectangle (3.1,3.25); \node[lblc] at (2.5,3.08) {\tiny HABduino};
  \draw[fill=green!40!black!40] (3.3,2.92) rectangle (4.5,3.25); \node[lblc] at (3.9,3.08) {\tiny OBC};
  \draw[fill=sdgray!40] (1.9,3.35) rectangle (3.1,3.6); \node[lblc] at (2.5,3.47) {\tiny SPOT};
  % środek masy
  \draw[sdnavy] (2.5,3.25) circle (0.12); \fill[sdnavy] (2.5,3.25) -- ++(0.12,0) arc (0:90:0.12) -- cycle;
  \fill[sdnavy] (2.5,3.25) -- ++(-0.12,0) arc (180:270:0.12) -- cycle;
  \draw[dashed,sdnavy] (2.5,-0.9) -- (2.5,6.4);
  \node[lbl,anchor=west,text=sdnavy] at (5.4,3.55) {środek masy w osi zawieszenia};
  % płyta na pokrywie i zwalniacz
  \draw[fill=blue!50] (0.9,4.7) rectangle (4.1,4.82);
  \draw[fill=sdgray] (2.7,4.82) rectangle (3.4,5.2);
  \node[lbl,anchor=west] at (5.4,4.75) {\textbf{płyta montażowa} na pokrywie: zwalniacz\\serwomechaniczny, mocowanie pasów};
  \draw[thick] (2.85,5.2) -- (2.5,6.4);
  \node[lbl,anchor=west] at (2.7,6.2) {linka do spadochronu / odcinacza};
  % antena GPS
  \draw[fill=black] (0.3,4.82) rectangle (0.75,4.95);
  \node[lbl,anchor=south] at (0.55,5.0) {GPS};
  % pasy
  \draw[sdorange,line width=2pt] (0.15,-0.05) -- (0.15,4.72);
  \draw[sdorange,line width=2pt] (4.85,-0.05) -- (4.85,4.72);
  \node[lbl,anchor=east,text=sdorange] at (-0.15,1.5) {pasy\\z regulatorami};
  % anteny w dół
  \draw[thick] (1.2,0) -- (1.2,-0.9);
  \draw[thick] (3.8,0) -- (3.8,-0.8);
  \node[lbl,anchor=west] at (5.4,-0.5) {anteny nadajników wyprowadzone w dół};
  % termometr
  \draw[thick] (5,1.7) -- (5.6,1.7);
  \node[lbl,anchor=west] at (5.7,1.7) {czujnik temperatury na zewnątrz};
\end{tikzpicture}}

% ---------------------------------------------------------------- harmonogram misji
\newcommand{\rysHarmonogram}{%
\begin{tikzpicture}[font=\sffamily\scriptsize,x=1.05cm,y=0.62cm]
  % oś: tygodnie od -16 do +2
  \foreach \w/\l in {0/T$-$16 tyg.,4/T$-$12,8/T$-$8,12/T$-$4,14/T$-$2,16/T0,17/T+1}{
    \draw[sdgray!50] (\w*0.62,0.6) -- (\w*0.62,-9.4);
    \node[lblc,anchor=south] at (\w*0.62,0.6) {\l};}
  \def\pasek##1##2##3##4##5{\fill[##5] (##2*0.62,-##1+0.32) rectangle (##3*0.62,-##1-0.32);
    \node[lbl,anchor=east] at (-0.15,-##1) {##4};}
  \pasek{0}{0}{2}{Cel misji, wymagania ładunku}{sdnavy}
  \pasek{1}{1}{9}{Projekt i budowa eksperymentu}{sdnavy!70}
  \pasek{2}{2}{10}{Przegląd platformy, naprawy, zakupy}{sdnavy!50}
  \pasek{3}{8}{13}{Testy (zimno, zasięg, odcinacz, czas pracy)}{sdorange!80}
  \pasek{4}{12}{12.4}{Przegląd gotowości (FRR)}{sdorange}
  \pasek{5}{11}{14}{Zgłoszenie PAŻP (dwa wnioski), ubezpieczenie}{ForestGreen!70}
  \pasek{6}{14.3}{16}{Gaz, predykcje codziennie}{ForestGreen!50}
  \pasek{7}{15.5}{15.75}{Konfiguracja trackera (T$-$2 dni)}{sdorange}
  \pasek{8}{15.75}{16}{Weryfikacja, logistyka (T$-$1)}{sdorange!60}
  \pasek{9}{16}{17}{Start, odzysk, raport (do T+7 dni)}{red!70}
  \draw[red!70,thick] (16*0.62,0.6) -- (16*0.62,-9.4);
\end{tikzpicture}}
